% Shared threat model for the three privacy experiments.
\subsection{Adversary and observables}
\label{sec:eval-threat}

We evaluate against a \emph{passive on-chain observer}: an adversary who sees every
transaction and event our approach emits, but never the plaintext policy, the witness,
or any client-side timing. This is the strongest realistic attacker for a public
blockchain---it has unrestricted, perpetual read access to the ledger.

Concretely, for each policy execution (a commit transaction followed by an access
transaction) the adversary observes the following on-chain features:
%
\begin{itemize}
  \item \textbf{Gas:} \texttt{commitGasUsed}, \texttt{accessGasUsed}.
  \item \textbf{Calldata size:} \texttt{commitCalldataLength}. \texttt{accessCalldataLength}.
  \item \textbf{Proof size:} \texttt{proofSize}.
  \item \textbf{Event logs:} \texttt{commitLogsCount}, \texttt{commitLogsSize},
    \texttt{accessLogsCount}, \texttt{accessLogsSize}.
\end{itemize}
%
By design, calldata length, proof size, and log counts/sizes are content-independent
(constant across policies), so any residual signal can only flow through the gas
features in the modeled public transcript. The effective feature set the attacker can
exploit therefore reduces to the two gas measurements (\texttt{commitGasUsed}, \texttt{accessGasUsed}); the remaining observables are retained in the pipeline but carry no policy-dependent variation. The defenses (commitment-only interface, randomized access-gas padding) target exactly this channel. We do not model other potential side channels such as block
timing, transaction ordering, mempool or broadcast patterns, or off-chain TPS
communication; compromised TPS behavior, adaptive requester probing, and off-chain side
channels remain outside this evaluation model, matching the system model in
Section~\ref{sec:problem}.

We consider three increasingly fine-grained attacker tasks, one per privacy property:
(i)~recover the policy's \emph{content} (PCH, Section~\ref{sec:eval-pch}); (ii)~tell two
\emph{semantically related} policies apart (PI, Section~\ref{sec:eval-pi}); and
(iii)~recover the policy's \emph{structure or parameter magnitude} (SPP,
Section~\ref{sec:eval-spp}). The primary attacker is supervised and label-aware; a
separate unsupervised clustering attacker provides a label-free cross-check
(Section~\ref{sec:eval-method}).
